
%\chapter{Introdução}


\chapter{Apresentação Gráfica}

É apresentado um resumo das normas para trabalhos acadêmicos da UDESC.

\section{Margens}

As margens devem seguir os seguintes padrões:

\begin{itemize}
\item Superior: 3,0 cm
\item Inferior: 2,0 cm
\item Interna: 3,0 cm
\item Externa: 2,0 cm
\end{itemize}

Parágrafo: a 1,25 cm da margem esquerda.

\section{Formatação do Papel e Fonte}

Papel branco, formato A4 (\SI{21}{cm} $\times$ \SI{29,7}{cm}), gramatura 75 com texto em cor preta 

\textbf{Elementos pré-textuais}: iniciar no anverso da página, com exceção da ficha catalográfica que deve vir no verso da folha de rosto.

Os elementos textuais e pós-textuais devem ser digitados no anverso e verso das páginas, com exceção das páginas que derem início as seções primárias, que devem ser no anverso das páginas.


\textbf{Formato}: Arial ou Times New Roman.

\textbf{Tamanho 12}: elementos pré-textuais, textuais e pós-textuais

\textbf{Tamanho 10}: citações longas, notas de rodapé, legendas das ilustrações, tabelas e paginação.

Obs.: A fonte tamanho 12 pt deve ser usada em todo o trabalho, inclusive títulos das seções e
subseções.

\section{ESPAÇAMENTO E PARÁGRAFOS}

O texto do corpo do trabalho deve ser \textbf{alinhado à esquerda e à direita (justificado)}, digitado
em \textbf{espaço de 1,5 linha}, em \textbf{fonte tamanho 12 }(com exceção das referências e citações, tratados adiante), incluindo apêndices e anexos

Citações longas devem ter um recuo de 4,0 cm da margem esquerda, espaço simples e fonte
em tamanho 10.

Notas de rodapé, legendas e fontes de ilustrações e tabelas devem ser apresentadas em
espaço simples (1,0 linha), em fonte tamanho 10.

As referências no final do trabalho devem ser alinhadas somente à margem esquerda,
separadas uma das outras por espaço duplo; entre as linhas da mesma referência deve ser usado
espaço simples.

Os títulos principais (seção primária) devem ser separados do texto por 1 (um) espaço de 1,5
linha em branco. Títulos não numerados devem ser centralizados, títulos numerados devem ser
alinhados à esquerda. A numeração deve ser separada dos títulos ou subtítulos por um espaço de
caractere (sem ponto).

Os títulos principais devem ser alinhados pela margem superior da mancha, sendo
apresentados sempre em nova página. Os subtítulos (seções secundárias, terciárias, etc.) devem
aparecer na sequência do texto (sem iniciar nova página), separados do texto anterior e posterior
por 1 (um) espaço de 1,5 linha em branco.

\section{Paginação}

Todas as páginas do trabalho, a \textbf{partir da folha de rosto}, devem ser contadas sequencialmente,
mas nem todas são numeradas. A numeração é colocada a partir da \textbf{primeira página da parte textual (Introdução)}, em algarismos arábicos, fonte 10.

No verso da página (páginas pares) o número de página é inserido dentro da margem
esquerda superior. No anverso da página (páginas ímpares) o número é inserido na margem direita
superior. Dessa forma, o número da página estará sempre na margem externa da mancha gráfica.

\section{Figuras}

As ilustrações devem ser colocadas após sua citação no texto, deixando-se um espaço
(1,5) entre o texto e a figura. Após, o texto prossegue a um espaço (1,5).

\textbf{Identificação}: parte superior, precedida da palavra designativa, número de ordem de
ocorrência no texto, em algarismos arábicos, travessão e título. Espaçamento
entrelinhas simples e fonte 12. Se o título ocupar mais de uma linha, a segunda linha
deverá iniciar abaixo da primeira palavra do título.

\textbf{Fonte (identificador do responsável)}: na parte inferior da ilustração deve-se indicar a
fonte consultada (elemento obrigatório, mesmo que seja elaborado pelo autor).
Apresentar a fonte consultada precedida da palavra Fonte em letra maiúscula e
minúscula, dois pontos e a fonte, em tamanho 10 pt.

Todas as ilustrações deverão ser \textbf{centralizadas} em relação a margem.

Se o espaço da página não permitir, a ilustração deve aparecer na página seguinte,
enquanto o texto prossegue normalmente no restante da página anterior.

Na Fig.~\ref{fig:marca_udesc_01} pode ser observada uma figura de exemplo.

\begin{figure}[!htbp]
	\centering	
 	\caption{Marca UDESC.}
	\includegraphics[width=0.5\textwidth]{cor_horizontal_pdf.pdf}
	\fonte{Elaborado pelo autor, \the\year.}	
	\label{fig:marca_udesc_01}
\end{figure}

\begin{figure}[!htbp]
\flushleft 
	\caption{Outras marcas da UDESC.}
	\begin{subfigure}[b]{0.5\linewidth}
		\centering
		\includegraphics[width=0.5\textwidth]{cor_vertical_pdf.pdf}
		\caption{Tipo de marca 01}
	\end{subfigure}%
	\begin{subfigure}[b]{0.5\linewidth}
		\centering
		\includegraphics[width=0.85\textwidth]{cor_horizontal_ass_1_pdf.pdf}
		\caption{Marca com descrição da UDESC.}
	\end{subfigure}
	\fonte{Retirado de (UDESC, 2018).}	
\end{figure}

\section{Tabelas}

A largura das tabelas e quadros, não poderá ultrapassar a configuração das margens esquerda
e direita do texto. Deve-se ajustar o tamanho das tabelas e quadros ao conteúdo.

Devem estar configurados em \textbf{espaçamento simples}, com \textbf{fonte 10 ou 12 pt}, e seguir o mesmo
padrão em todo o trabalho.

Com relação a formatação, a tabela apresenta os seguintes elementos: título, cabeçalho,
conteúdo, fonte e, se necessário, nota(s) explicativa(s) (geral e/ou específica).

As informações inseridas em uma tabela devem ser divididas por linhas na horizontal, porém as \textbf{bordas laterais não podem ser fechadas}.

Na Tabela~\ref{tab:tabela_exemplo} pode ser observada uma tabela de exemplo.

\begin{table}[!htbp]
\centering
\small
\renewcommand{\arraystretch}{1.1}
\caption{Parâmetros de comunicação entre o CCS e o ESP8266.}%
\label{tab:tabela_exemplo}
\begin{tabular}{C{4cm} | C{4cm} | C{4cm}}
\toprule
 & \multicolumn{2}{c}{\textbf{Dispositivos}} \\
\cmidrule(lr{0.5ex}){2-3}
\multirow{-2}{*}[2pt]{\textbf{Parâmetros}}	& \textbf{CCS}	& \textbf{ESP} \\ 
\midrule
Endereço IP 			& 192.168.137.01 	& 192.168.137.212  		\\
Porta de recebimento 	& 9001 				& 9001  				\\
Porta de envio 			& - 				& 9001  				\\
Endereço MAC 			& - 				& 60:01:94:0E:7E:20  	\\
\bottomrule
\end{tabular}
\vspace{2mm}
\fonte{Retirado de (EICHSTÄDT, 2018).}
\end{table}


\section{Equações e Fórmulas }

As equações e as fórmulas, quando forem apresentadas na sequência normal do texto, devem ser representadas em linha.

\textbf{Exemplo:}

[...] o quadrado da hipotenusa é igual à soma dos quadrados dos catetos, logo: $a^2+b^2=c^2$ onde $a$ representa [...]

Ao longo do texto, quando o mesmo contiver diversas fórmulas e equações, estas devem ser identificadas com números sequenciais, colocados entre parênteses, na extremidade direita da linha, junto à margem.

\textbf{Exemplo:}
\begin{equation} \label{eq:circ}
A = \pi r^2
\end{equation}
\begin{equation} \label{eq:poli}
8xz^3-3x^2+2x^2z+\frac{1}{2}y^3z - \frac{1}{3}x
\end{equation}

A equação \eqref{eq:circ} comparada com a equação \eqref{eq:poli}...


\textbf{Exemplo:}

A matriz de transformação que leva o sistema de coordenadas naturais para um sistema de componentes simétricas, utiliza o operador de deslocamento no tempo $a=e^{j\frac{2\pi}{3}}$, que produz uma rotação de $\pm \SI{120}{\degree}$ nos fasores de saída.
\begin{equation}
\begin{bmatrix}
	\tilde{s}_{abc}^+ \\[0.5em]
	\tilde{s}_{abc}^- \\[0.5em]
	\tilde{s}_{abc}^0 
\end{bmatrix}
= \frac{1}{3} \cdot
\begin{bmatrix}
	1 & a   & a^2 	\\
	1 & a^2 & a 	\\
	1 & 1 	& 1 
\end{bmatrix}
\cdot
\begin{bmatrix}
	\tilde{s}_a \\[0.5em]
	\tilde{s}_b \\[0.5em]
	\tilde{s}_c 
\end{bmatrix} .
\end{equation}


\textbf{Exemplo:}

\begin{equation} \label{eq:valor_final}
G_{i,k} = \prod_{i=A}^{k} G_i \Rightarrow 
\begin{cases}
\, \left| G_{i,k}  \right| =\   \displaystyle	 \prod_{i=A}^{k}  \big|\, G_i \, \big| \\[1.5em]
\angle G_{i,k}   =  \displaystyle\sum_{i=A}^{k}  \angle\, G_i\ 
\end{cases}.
\end{equation}

\textbf{Exemplo:}

\begin{equation}
A_p = 	\bigg(\frac{L_f\cdot I_{ef}\cdot I_{pk}}{B_{max}\cdot  J_{max}\cdot k_w}\bigg)^{1{,}14} = \SI{1}{cm^4}
\end{equation}









\chapter{Introdução}

\section{Introdução}

\subsection{Introdução}

\subsubsection{Introdução}

\subsubsubsection{Introdução}

%\lipsum[10]

%\begin{figure}[!htbp]
%	\label{fig:marca_udesc_01}
%	\centering	
% 	\caption{Marca UDESC.}
%	\includegraphics[width=0.5\textwidth]{cor_horizontal_ass_1_pdf.pdf}
%	\fonte{Elaborado pelo autor, \the\year}	
%\end{figure}
%
%\lipsum[11-13]
%
%\begin{figure}[!htbp]
%\flushleft 
%	\caption{Outras marcas udesc udesc udesc udesc udesc udesc udesc udesc udesc udesc udesc udesc udesc udesc udesc udesc udesc udesc udesc udesc udesc udesc udesc udesc udesc .}
%	\begin{subfigure}[b]{0.5\linewidth}
%		\centering
%		\includegraphics[width=0.5\textwidth]{cor_vertical_pdf.pdf}
%		\caption{Tipo de marca 01}
%	\end{subfigure}%
%	\begin{subfigure}[b]{0.5\linewidth}
%		\centering
%		\includegraphics[width=0.5\textwidth]{cor_vertical_ass_pdf.pdf}
%		\caption{Marca com descrição da UDESC.}
%	\end{subfigure}
%	\fonte{Elaborado pelo autor, \the\year}	
%\end{figure}
%
%\lipsum[14]
